\section{Data Description}
\label{sec:data_description}

This data release comprises multimodal magnetic resonance imaging (MRI) from a longitudinal visual-system study, organised according to the Brain Imaging Data Structure (BIDS) version~1.9.0 (\texttt{DatasetType}: \texttt{raw}). The packaged tree contains \textbf{84} participants (\texttt{sub-001}--\texttt{sub-084}) and \textbf{135} imaging sessions (\texttt{ses-01}: 83; \texttt{ses-02}: 52). Fifty-one participants have both sessions, 32 have \texttt{ses-01} only, and one participant (\texttt{sub-043}) has \texttt{ses-02} only; unavailable visits are omitted rather than represented by empty folders. Cohort labels in \texttt{participants.tsv} are Control ($n=56$), Glaucoma ($n=19$), DataON ($n=7$; optic neuritis), and DataTON ($n=2$; traumatic optic neuropathy); sex is recorded as Female ($n=46$) and Male ($n=38$). Demographic counts by cohort and session are summarised in Table~\ref{tab:demographics}. Released imaging modalities include anatomical MRI, blood-oxygen-level-dependent (BOLD) functional MRI, spin-echo field maps, and diffusion-weighted MRI, together with task event tables, physiological recordings, and selected derivatives. The overall organisation of the released dataset is illustrated in Fig.~\ref{fig:dataset_overview}.

%% Participant demographics by cohort and session (age omitted).
%% Layout matches reports/participant_summary/Participant_Demographics_by_Session_no_age.png
%% Values match reports/participant_summary/Participant_Demographics_by_Session_no_age.*
\begin{table}[htbp]
\centering
\caption{Participant demographic characteristics by cohort and imaging session (age omitted). Each cohort has ses-01 and ses-02 columns. Values are $n$ or $n$ (\%) as indicated.}
\label{tab:demographics}
\footnotesize
\setlength{\tabcolsep}{4pt}
\renewcommand{\arraystretch}{1.28}
%% tabular (not tabularx) inside resizebox: table keeps the figure's
%% 11-column layout and single-line cells, then scales to \textwidth.
\noindent\resizebox{\textwidth}{!}{%
\begin{tabular}{@{}l*{10}{c}@{}}
\toprule
\textcolor{headergray}{\textbf{Characteristic}} &
\multicolumn{2}{c}{\textcolor{headergray}{\textbf{Control}}} &
\multicolumn{2}{c}{\textcolor{headergray}{\textbf{\shortstack[c]{Optic neuritis\\(ON)}}}} &
\multicolumn{2}{c}{\textcolor{headergray}{\textbf{\shortstack[c]{Traumatic optic\\neuropathy\\(TON)}}}} &
\multicolumn{2}{c}{\textcolor{headergray}{\textbf{Glaucoma}}} &
\multicolumn{2}{c@{}}{\textcolor{headergray}{\textbf{Total}}} \\
\cmidrule(lr){2-3} \cmidrule(lr){4-5} \cmidrule(lr){6-7} \cmidrule(lr){8-9} \cmidrule(l){10-11}
 & \textcolor{headergray}{\textbf{ses-01}} & \textcolor{headergray}{\textbf{ses-02}}
 & \textcolor{headergray}{\textbf{ses-01}} & \textcolor{headergray}{\textbf{ses-02}}
 & \textcolor{headergray}{\textbf{ses-01}} & \textcolor{headergray}{\textbf{ses-02}}
 & \textcolor{headergray}{\textbf{ses-01}} & \textcolor{headergray}{\textbf{ses-02}}
 & \textcolor{headergray}{\textbf{ses-01}} & \textcolor{headergray}{\textbf{ses-02}} \\
\midrule
Participants, $n$
  & 55 & 41
  & 7 & 1
  & 2 & 1
  & 19 & 9
  & 83 & 52 \\
\rowcolor{tablealt}
Female, $n$ (\%)
  & 35 (63.6\%) & 25 (61.0\%)
  & 6 (85.7\%) & 1 (100.0\%)
  & 0 (0.0\%) & 0 (0.0\%)
  & 5 (26.3\%) & 3 (33.3\%)
  & 46 (55.4\%) & 29 (55.8\%) \\
Male, $n$ (\%)
  & 20 (36.4\%) & 16 (39.0\%)
  & 1 (14.3\%) & 0 (0.0\%)
  & 2 (100.0\%) & 1 (100.0\%)
  & 14 (73.7\%) & 6 (66.7\%)
  & 37 (44.6\%) & 23 (44.2\%) \\
\bottomrule
\end{tabular}%
}
\end{table}


\begin{figure}[htbp]
    \centering
    \begin{adjustbox}{max width=\linewidth, max height=0.72\textheight, center}
        \input{figures/fig1_dataset_structure}%
    \end{adjustbox}
    \caption{Overview of the structure of the released multimodal MRI dataset. The dataset is organised according to the Brain Imaging Data Structure (BIDS), with anatomical, functional, diffusion, and field-map data organised within participant and session directories. Task events, physiological recordings, quality-control outputs, processing derivatives, analysis code, and supporting documentation are also provided.}
    \label{fig:dataset_overview}
\end{figure}

At the dataset root, \texttt{dataset\_description.json} records the dataset name, BIDS version, raw dataset type, licence statement, and conversion provenance (\texttt{GeneratedBy}: \texttt{neuro\_pipeline} version~2.1.0). Participant-level metadata are limited to \texttt{participant\_id}, \texttt{cohort}, and \texttt{sex} in \texttt{participants.tsv}, with column definitions in \texttt{participants.json}. Additional root files include \texttt{README.md}, \texttt{CHANGES}, and task metadata sidecars \texttt{task-movie.json}, \texttt{task-movie\_events.json}, and \texttt{task-fmri\_events.json}. A \texttt{code/} directory provides stimulus-related resources used to interpret functional events (including \texttt{presentStimParams.m}, \texttt{task-fmri\_condition\_dictionary.tsv}, \texttt{task-fmri\_condition\_dictionary.json}, and \texttt{code/task-movie/}). Supporting documentation under \texttt{docs/} covers inventory, events, physiology, de-identification, BIDS validation, and quality-control summaries. The file \texttt{.bidsignore} excludes \texttt{*\_TB1TFL*} and \texttt{*\_localizer*} series from BIDS validation while leaving those files present in the tree.

Each packaged session follows the layout \texttt{sub-<id>/ses-<id>/}, with modality folders \texttt{anat/}, \texttt{func/}, \texttt{dwi/}, and \texttt{fmap/} when the corresponding data were collected, and a session-level \texttt{sub-<id>\_ses-<id>\_scans.tsv} listing imaging files in that session (\texttt{acq\_time} is recorded as \texttt{n/a}). Imaging volumes are distributed as gzip-compressed NIfTI (\texttt{*.nii.gz}) with JSON sidecars. Anatomical data under \texttt{anat/} include T1-weighted volumes labelled \texttt{*\_T1w.nii.gz}. Sidecar \texttt{ProtocolName}/\texttt{SeriesDescription} distinguish standard MPRAGE series (\texttt{T1w\_MPR}; 264 volumes in the release) from white-matter-nulled MPRAGE (\texttt{WMn\_MPRAGE\_sagittal}; 118 volumes). Five T1w volumes were withdrawn from the packaged tree after visual QC as unusable (\texttt{docs/quality\_control/excluded\_t1w.md}). Three-dimensional FLAIR volumes (\texttt{*\_FLAIR.nii.gz}; 137 files) are also released. Turbo-flash B1-mapping series (\texttt{*\_TB1TFL.nii.gz}; 272 files) and localizer images are present under \texttt{anat/} but are listed in \texttt{.bidsignore}. Released T1w and FLAIR anatomicals in the BIDS \texttt{anat/} hierarchy are the defaced volumes used for sharing; a separate \texttt{derivatives/defacing/} tree is not duplicated inside this package.

Functional data are stored under \texttt{func/} as magnitude BOLD series (\texttt{*\_bold.nii.gz}) for four task labels: \texttt{task-rest}, \texttt{task-movie}, \texttt{task-fmri}, and \texttt{task-control}. Coverage is incomplete across sessions: the release contains 135 magnitude \texttt{task-rest} runs (132 sessions; typically one run per session), 536 \texttt{task-movie} runs (134 sessions; most often four runs), 553 \texttt{task-fmri} runs (134 sessions; most often four runs), and 401 \texttt{task-control} runs (135 sessions; most often three runs). Representative volume counts in archived magnitude series are 320 (\texttt{task-rest}), 210 (\texttt{task-movie}), 226 (\texttt{task-fmri}), and 20 (\texttt{task-control}); individual runs may differ. Matching single-band reference images (\texttt{*\_sbref.nii.gz}) and phase images (\texttt{part-phase}) are included for many functional acquisitions. Event tables (\texttt{*\_events.tsv}) accompany all 536 magnitude \texttt{task-movie} runs and 514 of the 553 magnitude \texttt{task-fmri} runs; \texttt{task-rest} and \texttt{task-control} have no event files by design. Movie events include \texttt{onset}, \texttt{duration}, \texttt{trial\_type}, \texttt{stim\_file}, \texttt{eye}, \texttt{movie\_label}, and \texttt{matlab\_run\_id}. Task-fMRI events provide block onsets and \texttt{trial\_type} labels (\texttt{baseline} or \texttt{stim-01}--\texttt{stim-12}), with condition definitions in the root and \texttt{code/} dictionaries. Copyrighted movie video files are not redistributed; clip identity is recorded in the event tables and \texttt{code/task-movie/}. Spin-echo EPI field maps under \texttt{fmap/} use opposing phase-encode entities \texttt{dir-AP} and \texttt{dir-PA} (\texttt{*\_epi.nii.gz}; 268 anterior--posterior and 267 posterior--anterior magnitude series in the release, present in all 135 sessions) and are intended for susceptibility distortion correction of EPI data. Field-map sidecars include phase-encoding direction and, where available, \texttt{IntendedFor} references to associated functional or diffusion series.

Diffusion MRI under \texttt{dwi/} is released as \texttt{*\_dwi.nii.gz} with matching JSON sidecars and FSL-style \texttt{*.bval}/\texttt{*.bvec} gradient tables (911 magnitude diffusion series across 131 sessions). Protocol families in the sidecars include gSlider SMS acquisitions (\texttt{gsld\_76dir\_b2000\_1mmiso\_AP}, multi-shell $b=0/1000/2000$~s/mm$^2$; and short posterior--anterior $b=0$ series such as \texttt{gsld\_75TE\_PA\_3b0}) and RESOLVE trace acquisitions (\texttt{resolve\_3scan\_trace\_tra\_p3\_160\_1.4iso\_AP} and \texttt{resolve\_3scan\_trace\_tra\_p3\_160\_1.4iso\_PA}, including TRACEW series; $b=0/1000$~s/mm$^2$ in representative tables). Phase-encoding direction is recorded in the JSON sidecars and varies across runs.

Physiological recordings, when conversion succeeded, are stored alongside the corresponding functional runs as BIDS physiology files \texttt{*\_recording-<channel>\_physio.tsv.gz} with JSON sidecars. Channels present in the release are \texttt{pulse} (1127 files), \texttt{respiratory} (1258), and \texttt{trigger} (1223), totalling 3608 physiology time-series files spanning 77 participants and 116 sessions; absence of physiology for a given run does not by itself indicate invalid BOLD data. Derivatives included under \texttt{derivatives/} are (i) \texttt{mriqc/}, containing native MRIQC image-quality metric JSON for all 84 participants (HTML reports are not packaged), and (ii) \texttt{tractoflow\_normalize\_dwi/}, containing TractoFlow \texttt{Normalize\_DWI} outputs for 66 participants as \texttt{sub-<id>/ses-01/dwi/sub-<id>\_ses-01\_space-dwi\_desc-normalize\_dwi.\{nii.gz,json,bval,bvec\}}. Visual diffusion QC materials associated with that derivative are provided under \texttt{docs/quality\_control/dwi/}.
